% Propietario conservado: /Users/ruben/Documents/New project/output/REVISION_ARGUMENTAL_APERTURAS_CIERRES_20260915/VII/delivery/MANUSCRIPT_VII_ES_EN_COMPLETE/source_es/manuscrito/80_observador_y_respuesta_relacional.tex
% Recuperación y articulación de los propietarios documentados en
% PROCEDENCIA_OBSERVADOR_RESPUESTA.md. Originales conservados íntegros.
% Antecedentes de integración: núcleo APP--TRIT--TPK, realización del
% transporte de borde, lectores areal/volumétrico y reloj energético.
% Las condiciones de representación se enuncian antes de cada uso.

\section{Observador, memoria y respuesta relacional}
\label{vii:vii:sec:observador-respuesta}

La operación de lectura actúa sobre el transporte que publica el estado
enriquecido APP--TRIT--TPK. Su dominio conserva ruta, hoja, orientación,
acarreo y memoria; el resultado accesible al observador es una aplicación
de ese dominio. Esta distinción permite describir conjuntamente tres
operaciones: la elección de un representante orientado del transporte,
la inscripción de su registro en un aparato y la respuesta local inducida
por la frontera global. Cada operación transmite una información distinta.
La construcción siguiente especifica sus mapas y las condiciones que
permiten componerlos.

\subsection{Transporte de borde y lectura intrínseca}
\label{vii:vii:sec:observador-ciclo}

La publicación finita del transporte proporciona un ciclo orientado
\(C_n\), con \(n\geq1\), vértices \(0,\ldots,n-1\) y aristas
\(e_j:j\longrightarrow j+1\pmod n\). Fijamos la realización reticular
\(L\) del registro, su espacio vectorial
\(W=L\otimes_{\mathbb Z}\mathbb R\), un producto interno positivo y
el representante primitivo no nulo \(v\in L\) de la clase residual
distinguida, en la cámara declarada. Las hojas suma y producto de APP
aportan las polarizaciones del transporte; la orientación del TRIT y la
actualización del TPK preceden a esta publicación lineal. El ciclo,
la realización reticular y la clase residual son los datos que esta
sección recibe del transporte construido anteriormente.

Una \(0\)-cocadena es una función sobre los vértices con valores en
\(W\); una \(1\)-cocadena asigna un vector a cada arista orientada.
El coborde actúa por
\[
 (d\Phi)(e_j)=\Phi(j+1)-\Phi(j).
\]
Denotamos por \(\mathbf1_E\) la cocadena escalar constante en las
aristas. Para un transporte \(\omega\in C^1(C_n;W)\), su publicación
acumulada \(\Gamma_\omega\) se obtiene por sumas parciales e
interpolación afín; su defecto de cierre es
\begin{equation}
 B(\omega)=\Gamma_\omega(1)-\Gamma_\omega(0)
 =\sum_{j=0}^{n-1}\omega(e_j).
 \label{vii:vii:eq:observador-defecto}
\end{equation}

\begin{theorem}[Descomposición exacta y modo armónico del ciclo]
\label{vii:vii:thm:observador-descomposicion}
Todo transporte tiene una única descomposición
\[
 \omega=d\Phi+u\mathbf1_E,\qquad \Phi(0)=0,
 \qquad u=\frac1n\sum_j\omega(e_j).
\]
En particular, \(B(d\Phi)=0\) y \(B(\omega)=nu\).
\end{theorem}

\begin{proof}
Sea \(y_j=\omega(e_j)-u\). La suma de los \(y_j\) es cero.
Definimos \(\Phi(k)=\sum_{j=0}^{k-1}y_j\); el valor para \(k=n\)
coincide con \(\Phi(0)=0\), de modo que la función está bien definida
en el ciclo, incluida su arista final. Se cumple \(d\Phi(e_j)=y_j\).
Si \(d\Phi+u\mathbf1_E=d\Psi+w\mathbf1_E\), la suma de aristas da
\(nu=nw\). Entonces \(d(\Phi-\Psi)=0\); la conectividad del ciclo
y la normalización en el vértice inicial implican \(\Phi=\Psi\).
La identidad \(\sum_j(\Phi(j+1)-\Phi(j))=0\) prueba directamente
la cancelación telescópica.
\end{proof}

\begin{definition}[Transporte admisible y lectura normalizada]
\label{vii:vii:def:observador-admisible}
El transporte es admisible cuando su modo armónico pertenece a
\(\mathbb Rv\). Sobre el espacio \(\mathcal T_{\rm adm}\) de esos
transportes definimos la lectura global y su densidad por fase:
\begin{equation}
 \mathfrak a_v(\omega)
 =\frac{\langle B(\omega),v\rangle}{\langle v,v\rangle},
 \qquad
 \mu_v(\omega)=\frac{\mathfrak a_v(\omega)}n.
 \label{vii:vii:eq:observador-lectura}
\end{equation}
Se tiene \(B(\omega)=\mathfrak a_v(\omega)v\).
\end{definition}

\begin{theorem}[Determinación de la lectura intrínseca]
\label{vii:vii:thm:observador-unicidad}
El cociente
\(\mathcal T_{\rm adm}/dC^0(C_n;W)\) es unidimensional, generado
por la clase de \(v\mathbf1_E\). La aplicación \(\mu_v\) es el
único funcional lineal sobre \(\mathcal T_{\rm adm}\) que satisface
\[
 F(\omega+d\Phi)=F(\omega),\qquad F(v\mathbf1_E)=1.
\]
\end{theorem}

\begin{proof}
La descomposición anterior escribe cada transporte admisible como
\(d\Phi+\lambda v\mathbf1_E\). Su clase módulo cocadenas exactas
queda determinada por \(\lambda\). La clase de \(v\mathbf1_E\)
es no nula porque su suma es \(nv\neq0\). Por linealidad e
invariancia, todo funcional de las características indicadas vale
\(F(\omega)=\lambda F(v\mathbf1_E)=\lambda\). La fórmula
\eqref{vii:vii:eq:observador-lectura} da precisamente \(\mu_v=\lambda\).
La lectura global es \(\mathfrak a_v=n\mu_v\).
\end{proof}

\begin{definition}[Observador de borde]
\label{vii:vii:def:observador-borde}
Fijada la polarización del transporte publicado, un observador de borde
es un par \(\mathcal O=(\varepsilon,\Psi)\), donde
\(\varepsilon\in\{+1,-1\}\) y \(\Psi\in C^0(C_n;W)\).
Su acción de lectura es
\begin{equation}
 \omega^{\mathcal O}=\varepsilon\omega+d\Psi.
 \label{vii:vii:eq:observador-accion}
\end{equation}
El signo fija la orientación y la cocadena fija la trivialización.
Dos elecciones de polarización se comparan mediante su transporte
publicado; la relación \eqref{vii:vii:eq:observador-accion} describe
exactamente los representantes que difieren por orientación y coborde.
\end{definition}

\begin{proposition}[Covarianza de la lectura]
\label{vii:vii:prop:observador-covarianza}
La acción conserva la admisibilidad y verifica
\[
 B(\omega^{\mathcal O})=\varepsilon B(\omega),\quad
 \mathfrak a_v(\omega^{\mathcal O})=\varepsilon\mathfrak a_v(\omega),
 \quad \mu_v(\omega^{\mathcal O})=\varepsilon\mu_v(\omega).
\]
Las trivializaciones con igual orientación publican el mismo cierre.
\end{proposition}

\begin{proof}
La suma del coborde \(d\Psi\) es cero. Por tanto, el modo armónico
se transforma en \(\varepsilon u\), que sigue perteneciendo a
\(\mathbb Rv\), y el defecto se transforma en \(\varepsilon B\).
Las dos igualdades escalares se obtienen por la proyección normalizada.
\end{proof}

La unidad de esta lectura es el representante residual \(v\), con su
normalización interna. La libertad de trivialización actúa sobre el
representante de la cocadena; la información de cierre reside en su
clase. La descripción permite comparar observadores manteniendo
separados el transporte genealógico, su publicación lineal y el escalar
que evalúa la clase residual.

\subsection{Identidad de borde e interior celular}
\label{vii:vii:sec:observador-stokes}

Sea \(K\) un complejo celular finito orientado de dimensión dos,
provisto de una \(2\)-cadena fundamental \(c_K\) tal que
\(\partial c_K=[C_n]\). Esta condición expresa que las incidencias
interiores se cancelan con sus multiplicidades y orientaciones.
Sea \(\Omega\in C^1(K;W)\) una extensión de \(\omega\) y
\(F_\Omega=d\Omega\) su curvatura discreta aditiva.

\begin{theorem}[Identidad holográfica celular de la lectura]
\label{vii:vii:thm:observador-stokes}
Con los datos anteriores,
\begin{equation}
 B(\omega)=\langle\Omega,\partial c_K\rangle
          =\langle d\Omega,c_K\rangle,
 \qquad
 \mathfrak a_v(\omega)
 =\frac{\langle\langle d\Omega,c_K\rangle,v\rangle}
        {\langle v,v\rangle}.
 \label{vii:vii:eq:observador-stokes}
\end{equation}
En particular, \(F_\Omega=0\) implica \(\mathfrak a_v(\omega)=0\).
Si el cierre es no nulo, toda extensión del transporte en ese complejo
tiene curvatura discreta no nula.
\end{theorem}

\begin{proof}
El emparejamiento entre cadenas y cocadenas satisface
\(\langle d\Omega,c_K\rangle=\langle\Omega,\partial c_K\rangle\).
Al expandirlo, cada arista interior aparece con la suma de los números
de incidencia de las celdas contiguas; la condición sobre
\(\partial c_K\) deja exactamente el ciclo orientado de borde.
La restricción de \(\Omega\) a ese ciclo es \(\omega\), cuya suma
es \(B(\omega)\). La proyección sobre \(v\) da la segunda fórmula
y la implicación de planitud.
\end{proof}

La igualdad compara dos evaluaciones de la misma cocadena extendida:
el defecto acumulado del borde y el coborde integrado sobre el interior.
Su curvatura toma valores en \(W\) y usa el coborde celular aditivo.
La realización por una conexión geométrica conserva, adicionalmente,
su propio espacio de coeficientes y su ley de transporte.

\subsection{Inscripción isométrica del registro en el aparato}
\label{vii:vii:sec:observador-aparato}

Sea \(\mathfrak M_n\) un conjunto finito no vacío de registros
publicados a nivel \(n\), con resultado, ruta, acarreo, hoja y memoria.
Su espacio de registros es
\(\mathcal H_n=\ell^2(\mathfrak M_n)\), con base \(|m\rangle\).
La actualización visible \(F_n\) tiene un levantamiento de registro
\[
 \widehat F_n(m)=(F_n(m),m).
\]
Tomamos como registros del nivel siguiente un conjunto que contiene esas
parejas. La segunda coordenada conserva el antecedente y hace inyectivo
el levantamiento. Esta construcción especifica la memoria retenida por
la inscripción, también cuando la actualización visible reúne valores.

En un espacio finito de aparato \(\mathcal K_{A,n}\), fijamos una
familia ortonormal \((|a_m\rangle)_{m\in\mathfrak M_n}\) y definimos
\[
 \iota_n:\mathcal H_n\longrightarrow\mathcal K_{A,n},
 \quad \iota_n|m\rangle=|a_m\rangle,
 \qquad
 V_n:\mathcal H_n\longrightarrow\mathcal H_{n+1},
 \quad V_n|m\rangle=|\widehat F_n(m)\rangle.
\]
Sea \(\mathcal P_n=\iota_n\mathcal H_n\) el subespacio de puntero.

\begin{theorem}[Compatibilidad entre actualización e inscripción]
\label{vii:vii:thm:observador-entrelazador}
Los mapas \(\iota_n\) y \(V_n\) son isometrías. Existe una única
isometría \(W_n^0:\mathcal P_n\longrightarrow\mathcal K_{A,n+1}\)
que verifica
\begin{equation}
 W_n^0\iota_n=\iota_{n+1}V_n.
 \label{vii:vii:eq:observador-cuadrado}
\end{equation}
Si \(\dim\mathcal K_{A,n+1}\geq\dim\mathcal K_{A,n}\),
se prolonga a una isometría de todo \(\mathcal K_{A,n}\).
Mediante complementos auxiliares que igualen las dimensiones, se obtiene
una prolongación unitaria entre los espacios ampliados.
\end{theorem}

\begin{proof}
Las imágenes de las bases por \(\iota_n\) son ortonormales; las
imágenes por \(V_n\) también lo son por la inyectividad de
\(\widehat F_n\). Para \(z=\iota_n\psi\) definimos
\(W_n^0z=\iota_{n+1}V_n\psi\). La representación de \(z\) es única,
y los tres mapas conservan su norma. La fórmula determina el mapa en
una base de \(\mathcal P_n\) y prueba el cuadrado.
Si la dimensión de llegada es suficiente, el complemento ortogonal de
\(W_n^0\mathcal P_n\) contiene una familia ortonormal del tamaño de
\(\mathcal P_n^\perp\); enviando una base de este último a aquella
se construye la prolongación. Finalmente se amplían ambos espacios a
una misma dimensión y se completan las dos familias prescritas a bases
ortonormales. La aplicación entre esas bases es unitaria.
\end{proof}

Cuando el registro incluye rótulos de un instrumento finito, esta misma
inscripción realiza su ambiente: si los operadores
\(K_{y,a}:\mathcal H_S\to\mathcal H'_S\) satisfacen
\(\sum_{y,a}K_{y,a}^*K_{y,a}=I\), la aplicación
\[
 V_{\rm inst}\psi=\sum_{y,a}K_{y,a}\psi\otimes|y,a\rangle
\]
es isométrica, pues su norma al cuadrado es
\(\sum_{y,a}\langle\psi,K_{y,a}^*K_{y,a}\psi\rangle=\|\psi\|^2\).
La inscripción \(\iota\) transporta cada rótulo al puntero que le
corresponde. El modelo fija así el registro de la premedición; un
dispositivo material requiere además su interacción y su calibración.

\subsection{Álgebra accesible y complemento del puntero}
\label{vii:vii:sec:observador-expectativa}

Definimos \(P_m=|a_m\rangle\langle a_m|\),
\(P=\sum_mP_m\) y \(P_\perp=I-P\). El álgebra del subespacio de
puntero es \(P\mathcal B(\mathcal K_{A,n})P\), cuya unidad es \(P\).
En ella, el lector diagonal es
\[
 \mathcal D_P(A)=\sum_mP_mAP_m.
\]

\begin{proposition}[Expectativa del registro y extensión unital]
\label{vii:vii:prop:observador-pinzamiento}
La aplicación \(\mathcal D_P\) es una expectativa condicional
completamente positiva, preservadora de traza e idempotente sobre el
álgebra del puntero, con \(\mathcal D_P(P)=P\). Su imagen es
\(\bigoplus_m\mathbb CP_m\).
Sobre el álgebra del aparato completo, la aplicación
\begin{equation}
 \mathcal E_P(A)=\sum_mP_mAP_m+P_\perp AP_\perp
 \label{vii:vii:eq:observador-complemento}
\end{equation}
es una expectativa condicional unital, completamente positiva y
preservadora de traza. Su imagen es
\[
 \left(\bigoplus_m\mathbb CP_m\right)
 \oplus\mathcal B(\mathcal P_n^\perp).
\]
Ambas aplicaciones son proyecciones ortogonales para el producto de
Hilbert--Schmidt. Sus espectros están contenidos en \(\{0,1\}\).
Cuando existen componentes eliminadas, la separación entre el
subespacio fijo y el núcleo es uno.
\end{proposition}

\begin{proof}
Los proyectores \(P_m\) son mutuamente ortogonales. En el aparato
completo, la familia que añade \(P_\perp\) suma \(I\). Cada sumando
\(A\mapsto P_mAP_m\) es completamente positivo: tras cualquier
ampliación matricial actúa por congruencia con \(I\otimes P_m\).
Lo mismo ocurre para \(P_\perp\). La suma conserva esa propiedad.
La ciclicidad de la traza y la suma de los proyectores prueban la
preservación de traza en cada dominio. La ortogonalidad anula los
términos cruzados al iterar el mapa, por lo que es idempotente.
Sus elementos fijos son precisamente los bloques anunciados y la
aplicación es bimodular respecto de esa subálgebra.
Además,
\(\operatorname{Tr}(A^*P_mBP_m)
 =\operatorname{Tr}((P_mAP_m)^*B)\),
de donde resulta la autoadjunción para Hilbert--Schmidt. Una proyección
autoadjunta tiene únicamente los valores espectrales cero y uno cuando
ambos subespacios son no nulos. Si sólo queda el subespacio fijo,
la aplicación es la identidad.
\end{proof}

La suma \(\sum_mP_mAP_m\), aplicada sin el complemento a todo el
aparato, envía \(I\) a \(P\). Su unidad natural pertenece al álgebra
del puntero. La fórmula \eqref{vii:vii:eq:observador-complemento} conserva
también el bloque auxiliar; una resolución ortonormal de ese bloque
proporciona, cuando se desea, la diagonalización del aparato entero.

Si \(U\) es una transformación unitaria del aparato y
\(P'_m=UP_mU^*\), entonces
\[
 \mathcal E_{P'}(UAU^*)=U\mathcal E_P(A)U^*.
\]
La igualdad se obtiene sustituyendo cada proyector en la suma.
Asimismo, para \(A\in\mathcal B(\mathcal P_n)\), el cuadrado
\eqref{vii:vii:eq:observador-cuadrado} implica
\[
 \mathcal D_{P_{n+1}}(W_n^0A(W_n^0)^*)
 =W_n^0\mathcal D_{P_n}(A)(W_n^0)^*.
\]
En efecto, cada unidad matricial \(|a_m\rangle\langle a_{m'}|\)
se transporta a la unidad correspondiente a los registros distintos
\(\widehat F_n(m),\widehat F_n(m')\); el lector diagonal conserva
exactamente los términos con \(m=m'\).

\subsection{Factorización de la respuesta por la frontera efectiva}
\label{vii:vii:sec:observador-factorizacion}

Sea \(\mathcal X_\infty^{\rm enr}\) el dominio de estados compatibles
enriquecidos y sea
\(b:\mathcal X_\infty^{\rm enr}\longrightarrow B_\infty\) su traza
global de frontera. Denotamos por \(B_{\rm ef}=\operatorname{im}b\)
la frontera efectivamente publicada. Para un espacio de respuestas
\(\mathcal V_v\), consideramos
\(I_v:\mathcal X_\infty^{\rm enr}\longrightarrow\mathcal V_v\).

\begin{theorem}[Criterio de respuesta relacional]
\label{vii:vii:thm:observador-factorizacion}
Existe una única aplicación \(\overline I_v:B_{\rm ef}\to\mathcal V_v\)
tal que \(I_v=\overline I_v\circ b\) si y sólo si
\begin{equation}
 b(x)=b(x')\quad\Longrightarrow\quad I_v(x)=I_v(x').
 \label{vii:vii:eq:observador-fibra}
\end{equation}
La unicidad se extiende a todo \(B_\infty\) cuando \(b\) es
sobreyectiva.
\end{theorem}

\begin{proof}
Una composición por \(b\) es constante en cada fibra. Recíprocamente,
para \(\beta\in B_{\rm ef}\), se define
\(\overline I_v(\beta)=I_v(x)\) con \(b(x)=\beta\).
La condición \eqref{vii:vii:eq:observador-fibra} garantiza independencia
del representante; todo valor de \(\overline I_v\) queda determinado
por algún estado de esa fibra. Si \(b\) es sobreyectiva,
\(B_{\rm ef}=B_\infty\). En otro caso, la ecuación de composición
determina únicamente los valores sobre \(B_{\rm ef}\).
\end{proof}

Este criterio es conjuntista. Para una factorización continua, se
conservan además las topologías y la condición de aplicación cociente
de \(b:\mathcal X_\infty^{\rm enr}\to B_{\rm ef}\); con ella, la
continuidad de \(I_v\) equivale a la de su factor. Esta precisión
separa la determinación de una lectura de sus propiedades de regularidad.

Sea \(\ell_v:\mathcal X_\infty^{\rm enr}\to\mathcal L_v\) la
restricción accesible a una vecindad local. El carácter global efectivo
se verifica mediante dos estados que satisfagan
\begin{equation}
 \ell_v(x)=\ell_v(x'),\qquad I_v(x)\neq I_v(x').
 \label{vii:vii:eq:observador-testigo-global}
\end{equation}
Si \(I_v\) factoriza por \(b\), estos dos estados tienen fronteras
distintas. Además, la respuesta no puede escribirse como función de
\(\ell_v\), pues esa función asignaría el mismo valor a ambos estados.
La diferencia de fronteras adquiere así un efecto operatorio comprobable.

\subsection{Pesos de ambivalencia y respuesta autoinercial}
\label{vii:vii:sec:observador-ambivalencia}

Para la fórmula general consideramos lectores positivos
\(A_\partial,V_\partial\) y un lector energético \(Q\)
definidos sobre el dominio efectivo de frontera declarado.
A continuación se explicita una realización sobre
\(B_{\rm av}\subseteq B_{\rm ef}\). La memoria permanece como coordenada
\(\mathsf M_\partial(\beta)\) de la frontera. La bisagra de
ambivalencia proporciona las dos orientaciones
\[
 r_{\rm av}=\frac{\pi_{\rm HMT}}{\varphi_{\rm HMT}^2},\qquad
 r_{\rm av}^{J}=\frac{\varphi_{\rm HMT}}{\pi_{\rm HMT}^2}.
\]
El superíndice \(J\) expresa intercambio del par generador.
\paragraph{Realización normalizada de los lectores de frontera.}
La incidencia modal APP--TRIT--TPK, desarrollada en la construcción
regional, proporciona el envolvente de caminos de matriz
\[
 F_{\rm av}=\begin{pmatrix}0&1\\1&1\end{pmatrix}.
\]
El vector positivo \(v=(1,\varphi_{\rm HMT})^{\mathsf T}\)
satisface \(F_{\rm av}v=\varphi_{\rm HMT}v\). Su transición
normalizada y sus pesos estacionarios son
\begin{equation}
 P_{ij}=\frac{(F_{\rm av})_{ij}v_j}{\varphi_{\rm HMT}v_i},
 \qquad
 P=\begin{pmatrix}0&1\\
 \varphi_{\rm HMT}^{-2}&\varphi_{\rm HMT}^{-1}\end{pmatrix},
 \qquad
 (\mu_{\rm ar},\mu_{\rm vol})
 =\frac{(1,\varphi_{\rm HMT}^{2})}{1+\varphi_{\rm HMT}^{2}}.
 \label{vii:vii:eq:lectores-parry}
\end{equation}
Las filas de \(P\) suman uno por
\(\varphi^{-2}+\varphi^{-1}=1\); la ecuación estacionaria
\(\mu_{\rm ar}=\mu_{\rm vol}\varphi^{-2}\), junto con la
normalización, determina los pesos escritos. Esta medida de Parry
pondera los caminos del envolvente; la frecuencia cronológica
\((1/3,2/3)\) cuenta los instantes del calendario modal y conserva
su función distinta.

La realización de amplitud común de memoria se define sobre el sector
\(\mathcal X_{\rm av}\subset\mathcal X_\infty^{\rm enr}\) en el que
los lectores de las dos hojas tienen la factorización
\begin{equation}
 \mathcal L_{\rm vol}(x)=\mu_{\rm vol}\mathcal M(x),\qquad
 \mathcal L_{\rm ar}(x)=\pi_{\rm HMT}\mu_{\rm ar}\mathcal M(x),
 \qquad \mathcal M(x)>0.
 \label{vii:vii:eq:lectores-amplitud-comun}
\end{equation}
Ambos valores pertenecen a la misma recta positiva de amplitud de
memoria. La publicación areal lleva la marca regional
\(\pi_{\rm HMT}\) del retorno. La lectura adimensional se obtiene
dividiendo ambas secciones por \(\mathcal L_{\rm vol}(x)\).
De este modo las sumas de los pesos se efectúan entre magnitudes del
mismo tipo; una eventual realización como áreas y volúmenes físicos
se compone después con sus respectivos transductores dimensionales.

\begin{proposition}[Descenso de los lectores normalizados y del reloj energético]
\label{vii:vii:prop:lectores-frontera-explicitos}
Sobre \(B_{\rm av}=b(\mathcal X_{\rm av})\), las reglas
\begin{equation}
 A_\partial(b(x))
 =\frac{\mathcal L_{\rm ar}(x)}{\mathcal L_{\rm vol}(x)}
 =r_{\rm av},\qquad V_\partial(b(x))=1
 \label{vii:vii:eq:lectores-normalizados}
\end{equation}
definen dos lectores positivos independientes del representante.
Fijada una orientación \(a\) de la sección de acción, la energía de
decisión de \eqref{vii:v:ant:eq:ii-kb-aut-energia} determina sobre estados
\begin{equation}
 q_a(x)=\frac{h_a(x)}{108t_0(x)}
       =\frac{\hbar_a(x)}{t_P(x)}\in\mathcal L_E^+,
 \qquad t_P=\frac{54}{\pi_{\rm HMT}}t_0.
 \label{vii:vii:eq:lector-q-accion}
\end{equation}
Escribimos
\[
 b_{\rm av}:=b|_{\mathcal X_{\rm av}}:
 \mathcal X_{\rm av}\longrightarrow B_{\rm av}.
\]
Existe un único lector \(Q:B_{\rm av}\to\mathcal L_E^+\) tal que
\(q_a|_{\mathcal X_{\rm av}}=Q\circ b_{\rm av}\)
si y sólo si, para cualesquiera \(x,x'\in\mathcal X_{\rm av}\),
\begin{equation}
 b(x)=b(x')\ \Longrightarrow\
 \frac{h_a(x)}{t_0(x)}=\frac{h_a(x')}{t_0(x')}.
 \label{vii:vii:eq:descenso-reloj}
\end{equation}
En particular, una carta de acción y duración fijas satisface esta
condición y publica \(Q=h_a/(108t_0)\).
\end{proposition}
\begin{proof}
El cociente de \eqref{vii:vii:eq:lectores-amplitud-comun} cancela la
amplitud común y vale
\(\pi_{\rm HMT}\mu_{\rm ar}/\mu_{\rm vol}
=\pi_{\rm HMT}/\varphi_{\rm HMT}^2\).
El resultado es el mismo para todos los representantes de una fibra
de \(b_{\rm av}\) dentro de \(\mathcal X_{\rm av}\) y es positivo.
Dividir ambos lectores por el mismo factor positivo
conserva los cocientes ponderados en los que intervienen.
Para el reloj, el criterio de factorización del
teorema~\ref{vii:vii:thm:observador-factorizacion}, aplicado a \(q_a\),
equivale exactamente a \eqref{vii:vii:eq:descenso-reloj}. Las secciones
fijas lo satisfacen; también lo satisfacen secciones variables cuyo
cociente descienda por \(b_{\rm av}\). La unicidad resulta de la definición
de \(B_{\rm av}\) como imagen efectiva.
\end{proof}

El valor del lector \(Q\) pertenece a la recta de energía
\(\mathcal L_E=\mathcal L_{\rm act}\otimes\mathcal L_T^*\).
Con \(E_1,E_2\) en esa misma recta,
\(E_1E_2/Q^2\) es adimensional e independiente de la base energética
positiva. El símbolo \(Q\) de esta sección designa exclusivamente
esa energía de reloj; la cantidad de calor de la realización de
Landauer pertenece al balance térmico correspondiente.
La factorización de amplitud común especifica una realización
concreta de los lectores generales. Cada otro dominio conserva sus
aplicaciones de frontera y su propia comprobación del descenso.
La memoria completa sigue en la historia enriquecida y en su
publicación de frontera, aunque se cancele en el cociente normalizado.


Definimos
\[
 w_A(\beta)=
 \frac{A_\partial(\beta)r_{\rm av}}
 {A_\partial(\beta)r_{\rm av}+V_\partial(\beta)r_{\rm av}^{J}},
 \qquad w_V(\beta)=1-w_A(\beta).
\]
En la realización \eqref{vii:vii:eq:lectores-normalizados}, los pesos
admiten la evaluación exacta
\begin{equation}
 w_A=\frac{r_{\rm av}^{\,2}}{r_{\rm av}^{\,2}+r_{\rm av}^{J}}
     =\frac{\pi_{\rm HMT}^{4}}
            {\pi_{\rm HMT}^{4}+\varphi_{\rm HMT}^{5}},\qquad
 w_V=\frac{\varphi_{\rm HMT}^{5}}
            {\pi_{\rm HMT}^{4}+\varphi_{\rm HMT}^{5}}.
 \label{vii:vii:eq:pesos-amplitud-comun}
\end{equation}
La sustitución conserva los dos factores de la definición: el lector
areal normalizado y su ponderación posterior por \(r_{\rm av}\).
La generación y el marcado regional del retorno se realizan una vez.
Multiplicar numerador
y denominador por \(\varphi_{\rm HMT}^4\pi_{\rm HMT}^2\)
da la última expresión. Estos pesos permanecen constantes en el
sector de amplitud común. Con proyectores y energías fijos, una
variación de la respuesta puede proceder del lector \(Q\);
el testigo de dependencia global requiere además los estados con
igual restricción local y distinta respuesta de
\eqref{vii:vii:eq:observador-testigo-global}.


Sean \(P_A,P_V\) proyecciones ortogonales complementarias de un
espacio de respuesta \(\mathcal H_v\), y sean \(E_1,E_2>0\)
dos energías internas fijadas para la comparación. La respuesta es
\begin{equation}
 \overline I_v(\beta)=\frac{E_1E_2}{Q(\beta)^2}
       \bigl(w_A(\beta)P_A+w_V(\beta)P_V\bigr),
 \qquad I_v=\overline I_v\circ b.
 \label{vii:vii:eq:observador-respuesta-ponderada}
\end{equation}
También se admiten energías que sean lectores explícitos de
\(\beta\); en ese caso se conservan como tales en la fórmula.

\begin{proposition}[Positividad y covarianza de la respuesta]
\label{vii:vii:prop:observador-respuesta-positiva}
La respuesta \eqref{vii:vii:eq:observador-respuesta-ponderada} es
autoadjunta, positiva definida y constante en las fibras de \(b\).
Sus valores propios en los sectores no nulos son
\[
 \lambda_A=E_1E_2Q^{-2}w_A,\qquad
 \lambda_V=E_1E_2Q^{-2}w_V.
\]
El intercambio completo
\((A_\partial,r_{\rm av},P_A)
 \leftrightarrow(V_\partial,r_{\rm av}^{J},P_V)\)
permuta los dos sumandos y conserva el operador. Un transporte
unitario del espacio de respuesta conjuga ese operador y conserva
su espectro con multiplicidades.
\end{proposition}

\begin{proof}
La positividad de los lectores da \(0<w_A,w_V<1\). Para
\(z=z_A+z_V\), con \(z_A=P_Az\), \(z_V=P_Vz\),
\[
 \langle z,\overline I_vz\rangle
 =\frac{E_1E_2}{Q^2}
   (w_A\|z_A\|^2+w_V\|z_V\|^2)>0
\]
si \(z\neq0\). La descomposición ortogonal prueba las expresiones
espectrales. La dependencia exclusiva de \(\beta\) da constancia
en las fibras y la factorización. El intercambio de las ternas
intercambia numeradores, pesos y proyectores a la vez, por lo que
conserva la suma. Para un transporte unitario \(U\), la sustitución
\(P_A\mapsto UP_AU^*\), \(P_V\mapsto UP_VU^*\) produce
\(U\overline I_vU^*\), con el mismo espectro.
\end{proof}

Un intercambio implementado dentro de un mismo espacio por una unitaria
que envíe \(P_A\) a \(P_V\) requiere que ambos sectores tengan
igual dimensión. El intercambio de las dos ternas como rótulos de la
fórmula es válido con las dimensiones propias de cada sector.

La condición \eqref{vii:vii:eq:observador-testigo-global} exige comparar
las respuestas. Por ejemplo, multiplicar simultáneamente
\(A_\partial,V_\partial\) por un mismo factor positivo conserva
ambos pesos; si también se conservan \(E_1,E_2,Q\), la respuesta
permanece idéntica. En cambio, con iguales pesos y energías, sustituir
\(Q\) por \(qQ\), \(q>0\), multiplica la respuesta por \(q^{-2}\).
Si dos estados con igual restricción local realizan ese cambio con
\(q\neq1\), proporcionan el testigo global. Estas son condiciones
operatorias de separación; la existencia de los estados se comprueba
en el dominio de frontera que se esté realizando.

\subsection{Conexión del estado conjunto y lectura de aceleración}
\label{vii:vii:sec:observador-autoinercia}

En una realización diferenciable del estado conjunto
\(\Gamma=(x,\mathcal O,m)\), sean \(N\) su espacio de realización,
\(\nabla^{\rm tot}\) su conexión y \(\pi_S:N\to N_S\) una
proyección diferenciable hacia el subsistema, dotado de conexión
\(\nabla^S\). La conexión y la proyección conservan la procedencia
de los transportes y lectores que realizan. El marco observado se
obtiene mediante
\[
 \mathfrak F_{\mathcal O}(\Gamma)
 =\operatorname{Res}_{\mathcal O}(\Gamma,\nabla^{\rm tot}).
\]
Para una trayectoria dos veces diferenciable \(\Gamma(t)\),
escribimos \(\gamma_S=\pi_S\circ\Gamma\).

\begin{definition}[Autoinercia y defecto de proyección]
\label{vii:vii:def:observador-autoinercia}
Una trayectoria es autoinercial respecto de la conexión declarada
cuando \(\nabla^{\rm tot}_{\dot\Gamma}\dot\Gamma=0\).
El defecto de su proyección es
\begin{equation}
 \mathcal K_{\pi_S}(\dot\Gamma,\dot\Gamma)
 =\nabla^S_{\dot\gamma_S}\dot\gamma_S
  -d\pi_S\bigl(\nabla^{\rm tot}_{\dot\Gamma}\dot\Gamma\bigr).
 \label{vii:vii:eq:observador-defecto-proyeccion}
\end{equation}
\end{definition}

\begin{proposition}[Aceleración del subsistema]
\label{vii:vii:prop:observador-aceleracion}
Para una trayectoria autoinercial,
\[
 a_S=\nabla^S_{\dot\gamma_S}\dot\gamma_S
     =\mathcal K_{\pi_S}(\dot\Gamma,\dot\Gamma).
\]
En coordenadas \(x^B\) de \(N\) y \(y^a\) de \(N_S\), el
defecto es cuadrático en la velocidad, con componentes
\begin{equation}
 (\mathcal K_{\pi_S})^a{}_{BC}
 =\partial_B\partial_C\pi_S^a
  +(\Gamma^S)^a{}_{bc}\,
        \partial_B\pi_S^b\partial_C\pi_S^c
  -\partial_D\pi_S^a(\Gamma^{\rm tot})^D{}_{BC}.
 \label{vii:vii:eq:observador-defecto-coordenadas}
\end{equation}
\end{proposition}

\begin{proof}
La regla de la cadena da
\(\ddot\gamma_S^a=\partial_B\pi_S^a\ddot\Gamma^B
 +\partial_B\partial_C\pi_S^a\dot\Gamma^B\dot\Gamma^C\).
Sumando la contribución de \(\nabla^S\) y restando la imagen de
la aceleración total, los términos que contienen \(\ddot\Gamma\)
se cancelan. Los restantes son
\eqref{vii:vii:eq:observador-defecto-coordenadas} contraídos con
\(\dot\Gamma^B\dot\Gamma^C\). Al anularse la aceleración total
resulta la igualdad anunciada.
\end{proof}

La respuesta geométrica factoriza por \(b\) cuando la conexión,
la proyección y el dato cinemático utilizado en la evaluación se
determinan por \(\beta=b(x)\), o cuando esos datos cinemáticos se
mantienen fijados en la comparación. Con esas condiciones,
\eqref{vii:vii:eq:observador-defecto-coordenadas} da el mismo resultado
para cualesquiera representantes de una fibra; el
teorema~\ref{vii:vii:thm:observador-factorizacion} construye su factor
único sobre \(B_{\rm ef}\). Si varía el dato cinemático, éste
forma parte del dominio de la respuesta que se compara.

La formulación autoinercial describe una trayectoria geodésica del
estado conjunto y una aceleración efectiva de su subsistema debida
al defecto de proyección. La realización gravitatoria identifica
después esa conexión, el marco y la respuesta con sus observables
físicos. La covarianza del observador, la inscripción de memoria y
la factorización por la frontera conservan así funciones complementarias:
la primera determina la clase de lectura, la segunda retiene la
genealogía y la tercera especifica la dependencia relacional de la
respuesta.
